% ==============================================================================
% PLANTILLA MAESTRA GENERAL - UCNL
% ==============================================================================
% Uso:
% 1. Duplica este archivo y renombra la copia por materia, semana y actividad.
% 2. Actualiza documenttitle, documentsubtitle, documentsubject, coursename,
%    coursecode, universitydepartment y datos de la figura docente.
% 3. Lee la planeacion semanal y NO pegues las instrucciones completas dentro del
%    producto final. Usa la planeacion como lista de cumplimiento.
% 4. Revisa la bibliografia indicada, toma notas, registra citas y redacta con
%    comprension propia.
% 5. Identifica el producto academico solicitado y construyelo dentro del
%    documento siempre que sea posible.
% 6. Entrega en PDF con la nomenclatura indicada en el aula.
%
% Formato institucional recomendado:
% - Fuente tipo Helvetica, equivalente visual a Arial.
% - Interlineado 1.5.
% - Citas autor-anio con natbib: usar \citep{clave} o \citet{clave}.
% - Referencias en bibliografia BibTeX, formato APA 7 en cuerpo y referencias.
% - Caratula, resumen breve, desarrollo, producto solicitado, conclusion y
%   bibliografia.
% - Redaccion academica clara, coherente, cohesionada, formal y sin faltas.
% - La portada puede llevar una marca de agua institucional discreta; ajustar
%   las variables coverwatermark... si se requiere apagarla o cambiar intensidad.
%
% Criterios generales de cumplimiento:
% - Responder exactamente a la actividad solicitada en la planeacion.
% - Identificar conceptos, elementos, criterios o categorias de forma precisa.
% - Analizar, interpretar y tomar postura; no solo copiar fuentes.
% - Organizar la informacion por jerarquia, secuencia y criterios claros.
% - Incluir conclusion personal cuando la actividad lo permita o lo solicite.
% - Usar citas y referencias en APA 7.
% - Evitar copiar las indicaciones de la actividad en el cuerpo final.
%
% Regla para productos visuales:
% - Si la actividad pide tabla, cuadro comparativo, esquema, mapa conceptual,
%   mapa mental, linea de tiempo, diagrama, matriz, lista de verificacion, red
%   conceptual, SQA, organizador grafico o producto similar, realizarlo
%   preferentemente con LaTeX/TikZ para mantener nitidez, evidencia editable y
%   diseno controlado.
% - Si la consigna exige Canva, Miro u otra herramienta externa, insertar una
%   captura nitida y conservar la explicacion dentro del documento.
% - Para tablas extensas, usar pagina limpia, sin encabezados ni pies, landscape,
%   letra pequena y restaurar el estilo al terminar.
%
% Criterios de redaccion personal:
% - No escribir para llenar espacio; escribir para sostener una idea.
% - Convertir informacion en criterio: que significa, como funciona, que problema
%   resuelve, que limite presenta y que consecuencia produce.
% - Estructura mental recomendada: problema -> sistema -> consecuencia -> postura.
% - Usar pensamiento de diagnostico: entrada -> proceso -> salida.
% - Cada parrafo debe tener funcion: presentar, explicar, comparar, valorar o
%   concluir.
% - Usar primera persona con moderacion academica: "Considero que",
%   "Desde esta perspectiva", "La posicion que sostengo es". Evitar "yo creo".
%
% Notas de enfoque personal segun enfoque_personal.txt:
% - La voz debe ser academica, tecnica, directa, estrategica y funcional.
% - El texto no debe sonar como estudiante pasivo que recopila: debe sonar como
%   alguien que interpreta, ordena, cruza datos y busca consecuencia practica.
% - Antes de redactar, identificar la funcion de cada seccion: abrir problema,
%   definir criterio, demostrar con evidencia, comparar, valorar o concluir.
% - Evitar frases genericas como "es importante porque si" o "desde siempre".
%   Sustituirlas por relaciones verificables: causa, efecto, limite, costo,
%   funcion, riesgo, implicacion o mejora.
% - Cada concepto debe pasar de definicion a analisis. Formula util:
%   concepto -> funcion -> riesgo si se aplica mal -> consecuencia -> postura.
% - La postura personal no debe ser opinion suelta. Debe tener tres piezas:
%   posicion, razon y efecto. Ejemplo: "Considero que X, porque Y; por ello Z".
% - Integrar experiencia o criterio propio sin perder formalidad: no narrar la
%   vida personal, sino usar mirada practica para explicar por que el tema
%   importa en educacion, derecho, instituciones, comunicacion o profesion.
%
% Estilo observado en actividades ya realizadas:
% - Abrir con una tesis concreta:
%   "El Derecho no solo ordena conductas: organiza poder, define limites y
%   distribuye consecuencias."
% - Formular el problema en negativo y positivo:
%   "El problema no es solo..., sino..."
% - Convertir el producto visual en argumento:
%   "El cuadro no solo recopila datos; reconstruye la funcion, el problema y el
%   punto que requiere correccion."
% - Cerrar con aprendizaje transferible:
%   "No basta con saber que dice la norma; tambien debe analizarse a quien
%   protege, a quien excluye y que razones permiten defenderla publicamente."
%
% Nivel cognitivo esperado:
% - Recordar: usar solo para definir conceptos basicos. No debe ser el nivel
%   dominante en una actividad ponderable.
% - Comprender: explicar con palabras propias y relacionar con el tema.
% - Aplicar: llevar el concepto a un caso, norma, texto, correo, dilema o hecho.
% - Analizar: separar elementos, detectar causas, efectos, tensiones y limites.
% - Evaluar: sostener una postura con criterios, fuentes y consecuencias.
% - Crear: elaborar producto propio: mapa, cuadro, linea de tiempo, ensayo,
%   propuesta, reformulacion, glosario, portafolio o solucion argumentada.
%
% Regla de nivel minimo:
% - Si la actividad vale puntos y pide producto, no quedarse en recordar o
%   comprender. Debe llegar por lo menos a aplicar y analizar; si pide conclusion
%   o postura, debe llegar a evaluar.
%
% Uso de las 100 tecnicas didacticas:
% - Esta plantilla conserva el manual "Tecnicas didacticas del aprendizaje" que
%   se incorporo en las plantillas de Filosofia del Derecho, Etica y Moral
%   juridica, y Redaccion en contextos virtuales.
% - Cuando la planeacion mencione una tecnica de la coleccion UnADM, identificar:
%   definicion de la tecnica, producto esperado, pasos de construccion, criterios
%   de revision y cierre reflexivo.
% - La tecnica didactica no debe verse como adorno. Debe mostrar seleccion,
%   jerarquia, relacion entre ideas y lectura propia.
% - Siempre que sea viable, construir la tecnica en LaTeX/TikZ: cuadro
%   comparativo, mapa conceptual, esquema, mapa mental, linea de tiempo, matriz,
%   red conceptual, diagrama de flujo, tabla SQA o lista de verificacion.
%
% Checklist antes de entregar:
% [ ] El titulo, materia, semana, docente y fecha estan actualizados.
% [ ] El objetivo coincide con la planeacion.
% [ ] El producto solicitado aparece visible, rotulado y completo.
% [ ] La introduccion plantea un problema o postura, no solo anuncia el tema.
% [ ] El desarrollo explica, conecta y analiza.
% [ ] Hay citas APA autor-anio dentro del texto.
% [ ] La bibliografia final coincide con las fuentes citadas.
% [ ] La conclusion responde que se comprendio y que postura se sostiene.
% [ ] El nivel cognitivo llega al menos a aplicacion y analisis.
% [ ] Si hay conclusion/postura, incluye posicion, razon y consecuencia.
% [ ] La tecnica didactica solicitada esta construida y explicada.
% [ ] No aparecen instrucciones copiadas de la planeacion dentro del producto.
% [ ] Ortografia, acentos, sintaxis y nombres propios revisados.
% [ ] Si se uso IA como apoyo, el texto fue comprendido, revisado y apropiado.
% ==============================================================================

% CREACION DEL DOCUMENTO
\documentclass[
	spanish,
	letterpaper, oneside
]{article}

% ==============================================================================
% INFORMACION DEL DOCUMENTO
% Actualizar en cada actividad.
% Ejemplos:
% \def\documenttitle {Cuadro comparativo de conceptos fundamentales}
% \def\documentsubtitle {Actividad 1 - Administracion I}
% \def\documentsubject {Actividad 1}
% ==============================================================================
\def\documenttitle {La comunicacion en la administracion}
\def\documentsubtitle {Cuestionario resuelto: comunicacion formal, informal y organizacional}
\def\documentsubject {Actividad 1 - Administracion I}

\def\documentauthor {Martin Jonathan de la Cruz}
\def\coursename {Administracion I}
\def\coursecode {ADM-I}

\def\universityname {Universidad Ciudadana de Nuevo Leon}
\def\universityfaculty {Programa academico UCNL}
\def\universitydepartment {Administracion I}
\def\universitydepartmentimage {departamentos/logo-ucnl}
\def\universitydepartmentimagecfg {height=1.57cm}
\def\universitylocation {Monterrey, Nuevo Leon}

% INTEGRANTES, PROFESORES Y FECHAS
\def\authortable {
	\begin{tabular}{ll}
		\textbf{Alumno:}
		& \begin{tabular}[t]{l}
			 Martin Jonathan de la Cruz \\
		\end{tabular} \\
		Matricula: & UCNL \\

		\textit{Figura docente:}
		& \begin{tabular}[t]{l}
			Nombre \\ Apellido
		\end{tabular} \\
		& \\
		\multicolumn{2}{l}{Fecha de realizacion: \today} \\
		\multicolumn{2}{l}{\universitylocation}
	\end{tabular}
}

% IMPORTACION DEL TEMPLATE BASE
\input{template}

% FORMATO GENERAL
\usepackage{helvet}
\renewcommand{\familydefault}{\sfdefault}

% TABLAS Y PRODUCTOS VISUALES
\usepackage{array}
\usepackage{longtable}
\usepackage{pdflscape}
\usepackage{tikz}
\usetikzlibrary{
	arrows.meta,
	positioning,
	calc,
	matrix,
	fit,
	shapes.geometric,
	shadows.blur
}

% CITAS EN FORMATO AUTOR-ANIO, COMPATIBLE CON APA 7 EN EL CUERPO DEL TEXTO
% Usar:
% - Cita parentetica: \citep{clave}
% - Cita narrativa: \citet{clave}
% - Varias fuentes: \citep{clave1,clave2}
\setcitestyle{authoryear,open={(},close={)}}

% ==============================================================================
% AYUDAS DE REDACCION
% Quitar o comentar \pendiente cuando el documento final este terminado.
% ==============================================================================
\newcommand{\pendiente}[1]{}

% ==============================================================================
% MARCA DE AGUA EN PORTADA
% - Se aplica solo a la portada porque usa \AddToShipoutPictureBG*.
% - Para apagarla en alguna actividad: \def\coverwatermarkenabled {false}
% - Usar opacidad baja para que no compita con titulo, datos ni logotipo.
% ==============================================================================
\def\coverwatermarkenabled {true}
\def\coverwatermarkimage {img/departamentos/logo-nuevo-leon.png}
\def\coverwatermarkopacity {0.16}
\def\coverwatermarkwidth {0.72\paperwidth}
\def\coverwatermarkxshift {0cm}
\def\coverwatermarkyshift {-0.35cm}

\newcommand{\insertcoverwatermark}{%
	\ifthenelse{\equal{\coverwatermarkenabled}{true}}{%
		\AddToShipoutPictureBG*{%
			\begin{tikzpicture}[remember picture,overlay]
				\node[
					opacity=\coverwatermarkopacity,
					inner sep=0pt,
					xshift=\coverwatermarkxshift,
					yshift=\coverwatermarkyshift
				] at (current page.center) {%
					\includegraphics[width=\coverwatermarkwidth]{\coverwatermarkimage}%
				};
			\end{tikzpicture}%
		}%
	}{}%
}

\begin{document}

% PORTADA
\insertcoverwatermark
\templatePortrait

% CONFIGURACION DE PAGINA Y ENCABEZADOS
\templatePagecfg
\onehalfspacing

% ==============================================================================
% RESUMEN / ABSTRACT
% Debe ser breve. Formula recomendada:
% Tema + objetivo + tecnica/producto + enfoque personal.
% ==============================================================================
\begin{abstractd}
	\textbf{Objetivo:} Identificar los conceptos centrales de la comunicacion organizacional y distinguir la funcion de la comunicacion formal, informal, verbal y no verbal dentro de la administracion.

	\textbf{Producto elaborado:} Cuestionario resuelto y organizado en tabla sobre tipos de comunicacion, transparencia, retroalimentacion, resolucion de conflictos y uso de herramientas de comunicacion en la administracion.

	\textbf{Enfoque de analisis:} El cuestionario se responde desde una perspectiva funcional, centrada en como la comunicacion ordena procesos, reduce ambiguedades y fortalece la coordinacion organizacional.

	\textbf{Nivel cognitivo:} Comprender y aplicar, porque cada reactivo exige reconocer conceptos clave y asociarlos con su uso correcto dentro del contexto administrativo.
\end{abstractd}

% TABLA DE CONTENIDOS - INDICE
\templateIndex

% CONFIGURACIONES FINALES
\templateFinalcfg

% ======================= INICIO DEL DOCUMENTO =======================

% ==============================================================================
% SECCION 1: INTRODUCCION
% Apertura recomendada:
% - Iniciar con una afirmacion clara, no con una frase generica.
% - Traducir la consigna a un problema academico, juridico, etico, comunicativo
%   o profesional, segun la materia.
% - Explicar por que importa.
% - Cerrar anticipando el enfoque personal.
% ==============================================================================
\section{Introduccion}

La comunicacion administrativa no solo transmite informacion: tambien ordena funciones, coordina tareas y reduce la incertidumbre dentro de una organizacion. Cuando los mensajes circulan sin estructura, sin canales definidos o sin registro, se incrementan los malentendidos y se debilita la capacidad de dirigir equipos con claridad.

En este desarrollo, el problema de analisis consiste en distinguir que tipo de comunicacion corresponde a cada situacion y por que la comunicacion formal sigue siendo indispensable para documentar decisiones, asignar responsabilidades, resolver conflictos y sostener procesos de cambio. Por ello, el cuestionario se responde desde una perspectiva funcional: cada reactivo se interpreta segun el efecto que produce en la coordinacion, la transparencia y la confianza organizacional.

% ==============================================================================
% SECCION 2: DESARROLLO / ANALISIS
% Adaptar los subtitulos segun la materia y la planeacion.
%
% Productos posibles:
% - Cuadro comparativo: criterios, evidencia, diferencia, valoracion.
% - Mapa conceptual: concepto central, conceptos secundarios, conectores claros.
% - Linea de tiempo: fecha/etapa, hecho, aporte, consecuencia.
% - Estudio de caso: situacion, elementos, problema, interpretacion, solucion.
% - Lista de verificacion: criterio observable, indicador, cumplimiento, mejora.
% - Glosario: concepto, definicion, funcion, ejemplo, fuente.
% - Reporte o ensayo: tesis, argumentos, evidencia, postura y cierre.
%
% Regla de parrafo:
% Entrada: concepto o problema.
% Proceso: explicacion, fuente y relacion.
% Salida: consecuencia, criterio o postura.
% ==============================================================================
% ENCABEZADO TEMATICO (no 'Desarrollo'): nombra el contenido, no el proceso.
\section{Conceptos fundamentales de la comunicacion organizacional}

El estudio de la comunicacion administrativa parte de reconocer como circula la informacion en una organizacion y por que esa circulacion impacta la eficiencia del trabajo. El reto no esta solo en memorizar definiciones, sino en identificar que canal, que practica o que herramienta resulta mas adecuada para comunicar decisiones, orientar personal, prevenir conflictos y mantener coherencia institucional \citep{robbins2017administracion}.

La \textbf{comunicacion verbal} se realiza mediante palabras habladas o escritas; la \textbf{comunicacion no verbal} se apoya en gestos, expresiones faciales, posturas y otros elementos que complementan o matizan el mensaje \citep{chiavenato2019introduccion}. En el plano organizacional, la \textbf{comunicacion formal} sigue normas, protocolos y canales establecidos, mientras que la \textbf{comunicacion informal} surge de manera espontanea y puede ser util para agilizar relaciones, aunque no sustituye los registros institucionales cuando se requiere claridad, evidencia o seguimiento \citep{daft2019teoria}.

Tambien es importante distinguir la direccion del flujo comunicativo. La \textbf{comunicacion vertical} descendente transmite instrucciones de niveles superiores a niveles inferiores; la \textbf{comunicacion horizontal} articula a personas del mismo nivel; y la \textbf{comunicacion diagonal} conecta distintas areas y jerarquias cuando el trabajo exige coordinacion transversal \citep{koontz2012administracion}. A ello se suman elementos como la \textbf{retroalimentacion constructiva}, la \textbf{comunicacion no violenta} y el uso de \textbf{plataformas digitales}, que deben entenderse como medios para mejorar la colaboracion y no como sustitutos del criterio administrativo \citep{robbins2017administracion}. Cada reactivo del cuestionario aplica estos conceptos a situaciones concretas de la practica administrativa, de modo que resolverlo consiste en reconocer el principio comunicativo correcto y no en adivinar entre opciones \citep{unadm100TecnicasDidacticas2023}.

% METADISCURSO (comentado): nivel cognitivo aplicado. No visible; guia editorial.
% El nivel cognitivo dominante de la actividad es comprender y aplicar: reconocer el
% significado de conceptos (comunicacion formal, informal, verbal, no verbal, vertical
% o diagonal) y trasladarlos a situaciones concretas para seleccionar la opcion correcta.

\subsection{Cuestionario resuelto sobre comunicacion organizacional}

\clearpage
\pagestyle{empty}
\begin{landscape}
\begingroup
\scriptsize
\setlength{\tabcolsep}{5pt}
\renewcommand{\arraystretch}{1.25}
\begin{longtable}{@{}>{\centering\arraybackslash}p{0.035\linewidth}>{\raggedright\arraybackslash}p{0.40\linewidth}>{\raggedright\arraybackslash}p{0.19\linewidth}>{\raggedright\arraybackslash}p{0.30\linewidth}@{}}
\caption{Cuestionario resuelto sobre comunicacion organizacional.}\label{tab:cuestionario-comunicacion}\\
\toprule
\textbf{No.} & \textbf{Pregunta} & \textbf{Respuesta correcta} & \textbf{Justificacion} \\
\midrule
\endfirsthead
\toprule
\textbf{No.} & \textbf{Pregunta} & \textbf{Respuesta correcta} & \textbf{Justificacion} \\
\midrule
\endhead
\bottomrule
\endfoot
\bottomrule
\endlastfoot
1 & Cual de los siguientes es un tipo de comunicacion que se lleva a cabo mediante palabras habladas o escritas? & \textbf{b. Comunicacion Verbal.} & Es correcta porque b. Comunicacion Verbal corresponde con precision al concepto evaluado; las demas opciones no describen adecuadamente esa relacion. \\
2 & Cual es una directriz efectiva para la comunicacion formal? & \textbf{b. uso de formatos estandar como memorandos y cartas formales.} & Es correcta porque b. uso de formatos estandar como memorandos y cartas formales corresponde con precision al concepto evaluado; las demas opciones no describen adecuadamente esa relacion. \\
3 & Cual es el objetivo principal de la comunicacion formal en cuanto a roles y responsabilidades? & \textbf{a. clarificar roles y responsabilidades.} & Es correcta porque a. clarificar roles y responsabilidades corresponde con precision al concepto evaluado; las demas opciones no describen adecuadamente esa relacion. \\
4 & Cual es una tecnica para gestionar conflictos en la comunicacion organizacional? & \textbf{a. Usar comunicacion no violenta.} & Es correcta porque a. Usar comunicacion no violenta corresponde con precision al concepto evaluado; las demas opciones no describen adecuadamente esa relacion. \\
5 & Como contribuye la comunicacion formal a la transparencia organizacional? & \textbf{c. proporcionando un registro documentado de decisiones y politicas.} & Es correcta porque c. proporcionando un registro documentado de decisiones y politicas corresponde con precision al concepto evaluado; las demas opciones no describen adecuadamente esa relacion. \\
6 & La transparencia en la comunicacion organizacional ayuda a: & \textbf{d. Construir confianza entre empleados y direccion.} & Es correcta porque d. Construir confianza entre empleados y direccion corresponde con precision al concepto evaluado; las demas opciones no describen adecuadamente esa relacion. \\
7 & Que tecnica es crucial para resolver conflictos de manera constructiva? & \textbf{c. Comunicacion No Violenta.} & Es correcta porque c. Comunicacion No Violenta corresponde con precision al concepto evaluado; las demas opciones no describen adecuadamente esa relacion. \\
8 & Por que es importante la comunicacion formal para la toma de decisiones? & \textbf{b. proporciona un registro documentado y datos precisos.} & Es correcta porque b. proporciona un registro documentado y datos precisos corresponde con precision al concepto evaluado; las demas opciones no describen adecuadamente esa relacion. \\
9 & Que practica de comunicacion formal permite dar seguimiento a acuerdos, responsabilidades y datos usados en la toma de decisiones? & \textbf{b. documentar acuerdos, decisiones y datos relevantes.} & Es correcta porque b. documentar acuerdos, decisiones y datos relevantes corresponde con precision al concepto evaluado; las demas opciones no describen adecuadamente esa relacion. \\
10 & Que tipo de comunicacion formal fluye de niveles superiores a niveles inferiores en una organizacion? & \textbf{b. comunicacion vertical (de arriba hacia abajo).} & Es correcta porque b. comunicacion vertical (de arriba hacia abajo) corresponde con precision al concepto evaluado; las demas opciones no describen adecuadamente esa relacion. \\
11 & Que estrategia de comunicacion incluye el uso de plataformas como Microsoft Teams o Slack? & \textbf{b. Plataformas de Comunicacion Interna.} & Es correcta porque b. Plataformas de Comunicacion Interna corresponde con precision al concepto evaluado; las demas opciones no describen adecuadamente esa relacion. \\
12 & Cual es uno de los objetivos de la comunicacion efectiva en la administracion segun la introduccion? & \textbf{c. Explorar transformaciones tecnologicas.} & Es correcta porque c. Explorar transformaciones tecnologicas corresponde con precision al concepto evaluado; las demas opciones no describen adecuadamente esa relacion. \\
13 & Que aspecto de la comunicacion formal contribuye a la profesionalidad en un entorno de trabajo? & \textbf{c. adherirse a normas y protocolos establecidos.} & Es correcta porque c. adherirse a normas y protocolos establecidos corresponde con precision al concepto evaluado; las demas opciones no describen adecuadamente esa relacion. \\
14 & Cual es una de las principales ventajas de la comunicacion formal en la administracion? & \textbf{d. promueve la transparencia y la responsabilidad.} & Es correcta porque d. promueve la transparencia y la responsabilidad corresponde con precision al concepto evaluado; las demas opciones no describen adecuadamente esa relacion. \\
15 & La comunicacion que sigue normas y protocolos establecidos dentro de una organizacion se llama: & \textbf{d. Comunicacion Formal.} & Es correcta porque d. Comunicacion Formal corresponde con precision al concepto evaluado; las demas opciones no describen adecuadamente esa relacion. \\
16 & Que barrera de comunicacion puede surgir debido a la sobrecarga de informacion? & \textbf{c. Barreras de Sobrecarga de Informacion.} & Es correcta porque c. Barreras de Sobrecarga de Informacion corresponde con precision al concepto evaluado; las demas opciones no describen adecuadamente esa relacion. \\
17 & Que tipo de comunicacion se realiza a traves de gestos y expresiones faciales? & \textbf{c. Comunicacion No Verbal.} & Es correcta porque c. Comunicacion No Verbal corresponde con precision al concepto evaluado; las demas opciones no describen adecuadamente esa relacion. \\
18 & Cual es uno de los objetivos de proporcionar retroalimentacion constructiva a los empleados? & \textbf{b. Reforzar el sentido de logro y motivar.} & Es correcta porque b. Reforzar el sentido de logro y motivar corresponde con precision al concepto evaluado; las demas opciones no describen adecuadamente esa relacion. \\
19 & Que tipo de comunicacion formal se produce entre diferentes niveles jerarquicos y departamentos no directamente relacionados? & \textbf{a. comunicacion diagonal.} & Es correcta porque a. comunicacion diagonal corresponde con precision al concepto evaluado; las demas opciones no describen adecuadamente esa relacion. \\
20 & Cual es una caracteristica fundamental de la comunicacion formal? & \textbf{b. uso de lenguaje tecnico o especifico.} & Es correcta porque b. uso de lenguaje tecnico o especifico corresponde con precision al concepto evaluado; las demas opciones no describen adecuadamente esa relacion. \\
21 & La comunicacion escrita es crucial para: & \textbf{a. Documentar informacion y crear registros permanentes.} & Es correcta porque a. Documentar informacion y crear registros permanentes corresponde con precision al concepto evaluado; las demas opciones no describen adecuadamente esa relacion. \\
22 & Que teoria enfatiza la importancia de la participacion activa en los procesos de comunicacion organizacional? & \textbf{a. Teoria de la Comunicacion Participativa.} & Es correcta porque a. Teoria de la Comunicacion Participativa corresponde con precision al concepto evaluado; las demas opciones no describen adecuadamente esa relacion. \\
23 & Cual es un desafio en la comunicacion interna y externa? & \textbf{a. Mejorar la comunicacion intercultural.} & Es correcta porque a. Mejorar la comunicacion intercultural corresponde con precision al concepto evaluado; las demas opciones no describen adecuadamente esa relacion. \\
24 & Cual es una barrera psicologica que puede afectar la comunicacion? & \textbf{a. Percepciones personales y prejuicios.} & Es correcta porque a. Percepciones personales y prejuicios corresponde con precision al concepto evaluado; las demas opciones no describen adecuadamente esa relacion. \\
25 & Que estrategia es importante para superar barreras temporales en la comunicacion? & \textbf{b. Implementar tecnicas de gestion del tiempo.} & Es correcta porque b. Implementar tecnicas de gestion del tiempo corresponde con precision al concepto evaluado; las demas opciones no describen adecuadamente esa relacion. \\
26 & Cual es un desafio relacionado con la comunicacion tecnologica en la administracion? & \textbf{c. Asegurar que las herramientas tecnologicas no se vuelvan obstaculos.} & Es correcta porque c. Asegurar que las herramientas tecnologicas no se vuelvan obstaculos corresponde con precision al concepto evaluado; las demas opciones no describen adecuadamente esa relacion. \\
27 & La comunicacion informal se caracteriza por: & \textbf{b. Ocurrir de manera espontanea y no seguir canales establecidos.} & Es correcta porque b. Ocurrir de manera espontanea y no seguir canales establecidos corresponde con precision al concepto evaluado; las demas opciones no describen adecuadamente esa relacion. \\
28 & Que documento se utiliza para comunicar cambios en las politicas internas de una empresa? & \textbf{b. memorando.} & Es correcta porque b. memorando corresponde con precision al concepto evaluado; las demas opciones no describen adecuadamente esa relacion. \\
29 & Como contribuye la comunicacion efectiva a la gestion del cambio organizacional? & \textbf{b. Asegurando que los empleados comprendan y acepten el cambio.} & Es correcta porque b. Asegurando que los empleados comprendan y acepten el cambio corresponde con precision al concepto evaluado; las demas opciones no describen adecuadamente esa relacion. \\
30 & Cual es una caracteristica de la comunicacion no verbal? & \textbf{d. Complementa la comunicacion verbal.} & Es correcta porque d. Complementa la comunicacion verbal corresponde con precision al concepto evaluado; las demas opciones no describen adecuadamente esa relacion. \\
31 & Como puede la comunicacion interna afectar la cultura organizacional? & \textbf{c. Al transmitir valores y mision.} & Es correcta porque c. Al transmitir valores y mision corresponde con precision al concepto evaluado; las demas opciones no describen adecuadamente esa relacion. \\
\end{longtable}
\endgroup
\end{landscape}
\pagestyle{fancy}
\clearpage

% METADISCURSO (comentado): tecnica didactica aplicada. No visible; guia editorial.
% Se aplico una tecnica de cuestionario resuelto con organizacion tabular para hacer
% visible la relacion entre cada reactivo y su respuesta correcta; ordena el contenido
% para demostrar comprension y control del tema \citep{unadm100TecnicasDidacticas2023}.

% ==============================================================================
% BLOQUE OPCIONAL: TABLA O CUADRO COMPARATIVO AMPLIO
% Usar cuando el producto necesite mucho ancho:
% - \clearpage para iniciar en pagina limpia.
% - \pagestyle{empty} para quitar encabezado y pie en el producto amplio.
% - landscape para aprovechar el ancho horizontal.
% - scriptsize, tabcolsep y arraystretch para controlar legibilidad.
% - Restaurar \pagestyle{fancy} al terminar.
%
% Si la tabla es corta, puede omitirse landscape y usarse longtable normal.
% ==============================================================================
% \clearpage
% \pagestyle{empty}
% \begin{landscape}
% \begingroup
% \scriptsize
% \setlength{\tabcolsep}{4pt}
% \renewcommand{\arraystretch}{1.35}
% \begin{longtable}{@{}p{0.18\linewidth}p{0.39\linewidth}p{0.39\linewidth}@{}}
% \caption{Titulo del cuadro comparativo}\label{tab:cuadro-comparativo}\\
% \toprule
% \textbf{Criterio} & \textbf{Elemento 1} & \textbf{Elemento 2} \\
% \midrule
% \endfirsthead
% \toprule
% \textbf{Criterio} & \textbf{Elemento 1} & \textbf{Elemento 2} \\
% \midrule
% \endhead
% Concepto & \textit{Refuerzo Ciclo A: Contenido sintetico con cita si aplica.} & \textit{Refuerzo Ciclo A: Contenido sintetico con cita si aplica.} \\
% Caracteristicas & \textit{Refuerzo Ciclo A: Contenido.} & \textit{Refuerzo Ciclo A: Contenido.} \\
% Funcion & \textit{Refuerzo Ciclo A: Contenido.} & \textit{Refuerzo Ciclo A: Contenido.} \\
% Valoracion critica & \textit{Refuerzo Ciclo A: Contenido.} & \textit{Refuerzo Ciclo A: Contenido.} \\
% \bottomrule
% \end{longtable}
% \endgroup
% \end{landscape}
% \pagestyle{fancy}

% ==============================================================================
% BLOQUE OPCIONAL: PRODUCTO VISUAL EN TIKZ
% Usar para mapa conceptual, mapa mental, esquema, flujo, proceso o diagrama.
% Cada flecha debe expresar relacion; no usar nodos sueltos sin conectores.
% ==============================================================================
% \begin{figure}[H]
% 	\centering
% 	\begin{tikzpicture}[
% 		node distance=1.3cm,
% 		box/.style={rectangle, rounded corners=2pt, draw=black!65, fill=black!4, align=center, minimum width=3.2cm, minimum height=0.9cm},
% 		arrow/.style={-{Latex[length=2.5mm]}, thick, draw=black!65}
% 	]
% 		\node[box] (centro) {Concepto central};
% 		\node[box, right=of centro] (relacion) {Concepto relacionado};
% 		\node[box, below=of relacion] (efecto) {Consecuencia};
%
% 		\draw[arrow] (centro) -- node[above]{implica} (relacion);
% 		\draw[arrow] (relacion) -- node[right]{produce} (efecto);
% 	\end{tikzpicture}
% 	\caption{Titulo del producto visual.}
% \end{figure}

% ==============================================================================
% SECCION FINAL: CONCLUSION
% El analisis propio y la postura personal se integran ORGANICAMENTE en la
% conclusion (no como secciones separadas). Inicia en pagina propia (\clearpage).
% ==============================================================================
\clearpage
\section{Conclusion}

El cuestionario permitio identificar los conceptos basicos de la comunicacion organizacional y relacionarlos con situaciones propias de la administracion: la comunicacion formal aporta registro, responsabilidad y claridad; la no verbal complementa el sentido del mensaje; y la retroalimentacion, la comunicacion no violenta y las herramientas digitales solo son utiles cuando se integran con criterio administrativo. La solidez de cada respuesta se evaluo con criterios claros: su correspondencia con el concepto correcto y su fundamento en la teoria administrativa consultada \citep{robbins2017administracion,chiavenato2019introduccion}, no la mera memoria.

Desde mi analisis, el valor del ejercicio esta en reconocer que la mayoria de los reactivos converge en una misma idea: la comunicacion efectiva depende de claridad, registro, oportunidad y adecuacion del canal, de modo que no puede verse como un complemento secundario de la gestion, sino como una condicion operativa para coordinar personas y procesos. Considero que la comunicacion formal ocupa un lugar estrategico en la administracion porque convierte los acuerdos en acciones verificables y evita que la operacion dependa solo de interpretaciones personales; por eso sostengo que administrar bien no solo exige planear y dirigir, sino comunicar con precision, oportunidad y criterio etico.\footnote{Declaracion de uso de inteligencia artificial: se emplearon herramientas de IA como apoyo para organizar y depurar la redaccion; el analisis conceptual, la seleccion de respuestas y la postura personal son propios, revisados y validados por el autor, sin que la IA sustituyera el criterio academico.}

% ==============================================================================
% MANUAL DE TECNICAS DIDACTICAS
% Dejar incluido en la plantilla maestra. En una entrega final, puede comentarse
% esta linea si la docente solo solicita el producto de la semana.
% ==============================================================================
% Referencia opcional desactivada: archivo externo no disponible en este repositorio.

% BIBLIOGRAFIA
% Las referencias van en pagina propia (\clearpage), separadas de la conclusion.
\clearpage
\bibliography{administracion-I}

% FIN DEL DOCUMENTO

\end{document}


